We know how the DPLL algorithm works in general terms, a mechanism where \textbf{unit resolution} and \textbf{search} are combined together. \vspace{7pt}

Now we move to the implementation of the algorithm. A direct implementation of the pseudocode seen so far would make a copy of the \textit{CNF} $X$ at every
recursive call (\textit{$X$ is the propositional formula}), potentially going out of memory very quickly. \vspace{7pt}

What really happens is that somehow the implementation of the algorithm mimics the DPLL approach described previously, using a main data structure are performed 
the following operations:
\begin{enumerate}
    \item Unit resolution.
    \item Case Analysis.
    \item Backtrack.
    \item Fail.
\end{enumerate} \vspace{7pt}

So, basically we mimic the four main operations listed before inside a particular data structure, usually called $M$. The basic idea is: keep track of a list $M$
of literals (\textit{boolean varibles or their negations}) that have been decided and derived during the execution of DPLL.

More formally, the algorithm behaves following some \textbf{rules}:
\begin{enumerate}
    \item $M$ starts empty.
    \item $M$ is extended when:
    \begin{enumerate}
        \item \textbf{UnitPropagate}, a literal is derived by unit resolution;
        \item \textbf{Decide}, a new case analysis starts.
    \end{enumerate}
    \item $M$ is repaired when a contradiction is found:
    \begin{enumerate}
        \item \textbf{Backtrack}, go back to the last decision, remove everything behind the last decision, negate the decision, and continue with a new decision;
        \item \textbf{Fail}, otherwise, when it is not possible to backtrack, fail.
    \end{enumerate}
\end{enumerate}

Additionally, it uses a little bit of formal definitions about propositional logic, which are:
\begin{itemize}
    \renewcommand{\labelitemi}{-}
    \item A literal $l$ holds in $M$ ($M \models l$) if and only if $l$ occurs in $M$.
    \item A clause $C$ yields contradiction ($M \models \neg C$) if and only if for every literal $l$ in $C$, we have, $M \models \neg l$. It simply says: if all the literals of $C$ are present in $M$ in their negated form, we cannot satisfy anymore the clause.
    \item $l$ is \textit{undefined} in $M$ if and only if neither $M \models l$ nor $M \models \neg l$, so $l$ is not in the list $M$.
    \item A decision of a literal, described as $l^d$ (\textit{there isn't anything about mathematics, it expresses when a decision about a boolean variable was taken}), originates from a decision in the DPLL algorithm.
\end{itemize}